\documentclass[11pt,a4paper]{article}
\usepackage[T1]{fontenc}
\usepackage[utf8]{inputenc}
\usepackage[margin=2.7cm]{geometry}
\usepackage{amsmath,amssymb}
\usepackage{booktabs}
\usepackage[expansion=false]{microtype}
\usepackage{enumitem}
\usepackage[hidelinks]{hyperref}

\setlist[itemize]{itemsep=2pt,topsep=3pt}
\newcommand{\Es}{E^{*}}
\newcommand{\kcal}{\,kcal/mol}

\title{ProbJanusDDG: A Protein Stability Change Predictor with Calibrated Uncertainty Intervals for Individual Mutations}
\author{}

\begin{document}
\maketitle
\vspace{-10mm}




\section{Abstract}
In recent years, numerous tools have been developed to predict the change in protein stability upon mutation ($\Delta\Delta G$). However, their accuracy has not significantly improved, likely due to a lack of new training data. Accuracy is not the only factor of practical importance when applying these methods: obtaining an uncertainty estimate on the predicted output is equally crucial. The lack of state-of-the-art methods that return a validated uncertainty alongside the point prediction of $\Delta\Delta G$ represents a critical gap to bridge to make these models more usable.

%In this paper, we present an extension of JanusDDG, a recent state-of-the-art model for predicting protein stability changes under mutations, designed to provide a validated confidence interval for every model prediction. 

In this paper, we present an extension of JanusDDG, a recent state-of-the-art model for predicting mutation-induced changes in protein stability, designed to associate each prediction with an uncertainty interval with empirically validated coverage.

To achieve this, we employed a Gaussian output head by replacing the single-neuron output layer of JanusDDG with a two-neuron layer parameterising the Gaussian mean $\mu$ and standard deviation $\sigma$. 
We conformally calibrated the model across the S2450 training folds, using quantiles of the prediction errors scaled by median-normalised $\sigma$ to determine interval widths. We then
evaluated interval coverage on the external test sets S669 and S461, both featuring low sequence identity to the training set.

%Combining $\sigma$ with conformal calibration across the S2450 training folds, we evaluated interval coverage on the external test sets S669 and S461, both featuring low sequence identity to the training set.

Model accuracy remains competitive with the state of the art ($r=0.54$ on S669, $0.70$ on S461), and prediction interval coverage transfers robustly to external benchmarks. The primary contribution of $\sigma$ is to redistribute interval width without increasing the mean interval width: a constant conformal interval is unnecessarily wide for mutations the model deems easy and overly narrow for those where it is uncertain. Stratifying S669 into three tiers (easy, medium, and hard mutations) by predicted $\sigma$, the constant interval achieves 97.6\%, 89.8\%, and 81.9\% coverage against the nominal 90\% target across all three groups; standardising by $\sigma$ yields 93.7\%, 87.8\%, and 86.4\%, effectively halving the gap between easy and hard tiers while largely eliminating the corresponding in-domain coverage imbalance.


%eliminating it entirely in-domain.


Having established model robustness, 
this work delivers a $\Delta\Delta G$ predictor whose accuracy is competitive with the state of the art and whose every prediction carries an interval with empirically verified coverage, tighter where the model is confident and wider where it is not.


\end{document}